% =================================================================
%  style/math-commands.tex
%  Математические пакеты, сокращения, нормы и скобки
% =================================================================

\usepackage{amsmath,amssymb,amsthm,mathtools,cancel}
\usepackage{mathrsfs}

% -----------------------------------------------------------------
% Числовые множества
% -----------------------------------------------------------------
\newcommand{\NN}{\mathbb{N}}
\newcommand{\ZZ}{\mathbb{Z}}
\newcommand{\QQ}{\mathbb{Q}}
\newcommand{\RR}{\mathbb{R}}

\providecommand{\R}{\mathbb{R}}
\providecommand{\C}{\mathbb{C}}
\providecommand{\N}{\mathbb{N}}

% -----------------------------------------------------------------
% Частые символы
% -----------------------------------------------------------------
\newcommand{\eps}{\varepsilon}
\newcommand{\vphi}{\varphi}
\newcommand{\st}{\;:\;}
\newcommand{\bigO}{\mathcal{O}}
\providecommand{\eps}{\varepsilon}

\DeclareMathOperator{\dist}{dist}

% -----------------------------------------------------------------
% Нормы, модули, скалярное произведение
%   \Lp{f}      -> ||f||_{L^p}
%   \Lq{f}      -> ||f||_{L^q}
%   \norm{f}{2} -> ||f||_{L^2}
%   \norma{x}   -> ||x||        (растягивающиеся скобки)
%   \abs{x}     -> |x|
%   \scal{x}{y} -> <x, y>
% -----------------------------------------------------------------
\newcommand{\Lp}[1]{\bigl\|#1\bigr\|_{L^p}}
\newcommand{\Lq}[1]{\bigl\|#1\bigr\|_{L^q}}
\newcommand{\norm}[2]{\bigl\|#1\bigr\|_{L^{#2}}}

\providecommand{\norma}[1]{\left\lVert #1 \right\rVert}
\providecommand{\abs}[1]{\left\lvert #1 \right\rvert}
\providecommand{\scal}[2]{\left\langle #1,#2 \right\rangle}

% -----------------------------------------------------------------
% Классические theorem-окружения (нумерация по section).
% В тексте используются цветные боксы из theorem-boxes.tex;
% эти оставлены как «тихий» вариант без рамки.
% -----------------------------------------------------------------
\theoremstyle{plain}
\newtheorem{theorem}{Теорема}[section]
\newtheorem{lemma}[theorem]{Лемма}
\newtheorem{proposition}[theorem]{Предложение}
\newtheorem{statement}[theorem]{Утверждение}

\theoremstyle{definition}
\newtheorem{definition}[theorem]{Определение}
\newtheorem{example}[theorem]{Пример}

\theoremstyle{remark}
\newtheorem{remark}[theorem]{Замечание}

% -----------------------------------------------------------------
% Доказательство: «Доказательство.» жирным, в конце — чёрный квадрат
% -----------------------------------------------------------------
\renewcommand\qedsymbol{$\blacksquare$}

\makeatletter
\renewenvironment{proof}[1][\proofname]{%
	\par\pushQED{\qed}%
	\normalfont \topsep6\p@\@plus6\p@\relax
	\trivlist
	\item\relax\textbf{#1.}\hspace\labelsep\ignorespaces
}{%
	\popQED\endtrivlist\@endpefalse
}
\makeatother
